\chapter*{Conclusiones}
\addcontentsline{toc}{chapter}{Conclusiones}
\markboth{Conclusiones}{Conclusiones}

\lettrine[lines=2]{E}{n} este trabajo se desarrolló un formalismo vectorial para
calcular el campo electromagnético en el plano focal de un sistema estigmático
simple, formado por una única superficie cartesiana que separa dos medios
homogéneos. El formalismo combina la geometría exacta de la superficie, la
refracción vectorial dada por los coeficientes de Fresnel y la propagación de
Fresnel hasta el plano focal, y produce el campo focal como transformada de
Fourier de una función pupila vectorial generalizada
$\mathbb E^{\mathscr G'}=\mathbb T\,\mathbb E_0/\ell_0$; bajo iluminación
azimutalmente simétrica esta transformada se reduce a transformadas de Hankel
de órdenes cero, uno y dos. A diferencia del tratamiento canónico de Richards
y Wolf~\cite{RichardsWolf1959}, formulado para sistemas aplanáticos
idealizados, aquí la superficie refractora y su transmisión de Fresnel entran
de forma explícita en la función pupila, punto a punto.

El resultado central es que un sistema geométricamente perfecto no está libre
de aberraciones. La condición estigmática elimina la aberración de camino
óptico, pero al cambiar de medio las condiciones de frontera de Maxwell
inducen sobre el campo una transformación vectorial anisótropa y distinta en
cada punto de la superficie. Sobre el plano transversal, esta transformación
es la de un diatenuador orientado radialmente en cada punto de la pupila,
cuyos modos propios son las polarizaciones radial y azimutal, moduladas
respectivamente por $t_p$ y $t_s$. Las combinaciones $t_\pm=t_p\pm t_s$
separan los dos efectos de la refracción: $t_+$ es una apodización común a
ambas componentes, mientras que la aberración vectorial, definida como la
desviación del campo focal respecto del límite escalar $t_s=t_p=1$ (en el
cual el formalismo recupera exactamente la teoría escalar de la difracción),
aparece siempre modulada por el factor $t_-$, no nulo y no uniforme para
cualquier interfaz entre medios de distinto índice. En analogía con la teoría
de aberraciones de frente de onda, la aberración queda así caracterizada
sobre la propia pupila por la función matricial $\mathbb T$.

Las manifestaciones observables dependen del estado de polarización
incidente. En la base cilíndrica la transmisión es diagonal: las
polarizaciones radial y azimutal no sufren mezcla, y bajo iluminación
azimutalmente simétrica la estructura de Hankel de orden uno anula el campo
transversal en el eje, produciendo un anillo con centro oscuro. En las bases
que no coinciden con los modos propios, $t_-$ acopla polarizaciones
ortogonales: en la circular genera un canal
$\ket{L}\!\leftrightarrow\!\ket{R}$ con fase azimutal $e^{\pm2i\varphi}$,
suprimido sobre el eje donde $t_-$ y $J_2$ se anulan simultáneamente; en la
cartesiana genera una componente cruzada de orden angular dos, visible como
cuatro lóbulos en la polarización ortogonal a la incidente, que sobre el
estado de polarización actúa como una rotación local de la orientación sin
introducir elipticidad. Distintas polarizaciones incidentes producen así
distintos patrones de intensidad focal, algo que la teoría escalar, cuya
predicción es única e independiente de la polarización, no puede describir.

El régimen vectorial modifica también la forma del punto focal y la imagen de
objetos extensos. Para incidencia lineal la interferencia entre los canales
$\mathscr H_0$ y $\mathscr H_2$ rompe la simetría circular del lóbulo
central: aparecen dos radios principales que satisfacen
$R_\parallel<R_{\mathrm{esc}}<R_\perp$, de modo que el punto se contrae a lo
largo de la polarización incidente, se expande en la dirección perpendicular
y rota rígidamente con ella. En el dioptrio de vidrio estudiado la
deformación es del orden del $3\,\%$ ($e=R_\perp/R_\parallel\approx1.06$) y
crece con la apertura y el contraste de índice hasta $e\approx1.14$ en un
dioptrio rápido de alto índice; con incidencia circular el punto permanece
circular y solo se ensancha levemente. La misma ruptura de isotropía aparece
en la imagen de una máscara litográfica formada por una lente de dos
superficies cartesianas conjugadas a $\mathrm{NA}=0.94$, libre de aberración
geométrica por construcción: las trazas alineadas con una iluminación
linealmente polarizada conservan la mayor parte del contraste
($C_H:0.98\to0.85$) mientras que las ortogonales se degradan severamente
($C_V:0.91\to0.67$), y la iluminación circular reparte la pérdida entre ambas
orientaciones sin eliminarla. En este régimen de apertura el mecanismo
dominante es la componente longitudinal $E_z$, determinada por la condición
de transversalidad, que transporta cerca de la mitad de la energía de la
imagen. Se trata del mismo fenómeno que obliga a la litografía de alta
apertura numérica a controlar la polarización de la
iluminación~\cite{flagello1996, ruoff2009}, presente aun cuando el sistema
está libre de aberración geométrica.

El formalismo admite extensiones naturales: la reflexión se obtiene
reemplazando $t_p,t_s$ por $r_p,r_s$; los sistemas de varias superficies se
tratan por composición de operadores, como ilustra el ejemplo litográfico; y
el dioptrio esférico en los puntos de Young--Weierstrass ofrece un caso de
validación de especial interés por su realizabilidad práctica. La formulación
se generaliza además a iluminaciones arbitrarias y a superficies dióptricas
no necesariamente estigmáticas, en cuyo caso la aberración de camino óptico
propia de la geometría se superpone a la aberración vectorial aquí
identificada. La verificación experimental de las firmas predichas, en
particular la componente cruzada de cuatro lóbulos y la deformación
anisótropa del punto focal, constituye la continuación natural de este
trabajo.
